\section{Version 2}

\subsection{Proposition}

A proposition is written R:RelT(x: LeftT, y: RightT) and it is equivalent to the following lambda expression:

\begin{equation}
\begin{split}
\lambda R:RelT.\lambda x:LeftT.\lambda y:RightT.(R\ x\ y)
\end{split}
\end{equation}

A proposition $P$ means that $\forall x \in \text{dom(LeftT)}$ and $\forall y \in \text{dom(RightT)}$, $P$ is a proof such that $x R y$ holds.
The domain of $P$ is the set of tuples $(R, x, y)$ such that $x R y$ holds.

A proposition $\overline{P}$ is a proof such that $x R y$ does not hold.
The domain of $\overline{P}$ is the empty set.

\subsection{Measures}

Instances of relations $aRb$ entering the system are called measures.

\subsection{Objective}

The objective of the system is to find the most compact set of hypotheses explaining the data, i.e., leading to a proof for each data.
